\documentclass[tikz,border=4pt]{standalone}
\usepackage{ctex}
\usepackage{amsmath}
\usetikzlibrary{arrows.meta}

% p8-05b: 两直线三种位置关系（交叉/平行/重合）对应方程组解的个数
\begin{document}
\begin{tikzpicture}[scale=0.75, thick]

  % ====== 左：相交（唯一解）======
  \begin{scope}[shift={(-4.5, 0)}]
    \draw[->] (-1.5,0) -- (1.5,0) node[right] {\tiny$x$};
    \draw[->] (0,-1.5) -- (0,1.8) node[above] {\tiny$y$};
    \node[below left] at (0,0) {\tiny$O$};
    \draw[blue, thick] (-1.2,-1.2) -- (1.2,1.2);
    \draw[red, thick] (-1.2,1.0) -- (1.2,-0.4);
    \fill (0.3, 0.3) circle (2.5pt);
    \node[above right, font=\scriptsize] at (0.3,0.3) {交点};
    \node[below, font=\scriptsize] at (0,-1.8) {唯一解（相交）};
  \end{scope}

  % ====== 中：平行（无解）======
  \begin{scope}[shift={(0, 0)}]
    \draw[->] (-1.5,0) -- (1.5,0) node[right] {\tiny$x$};
    \draw[->] (0,-1.5) -- (0,1.8) node[above] {\tiny$y$};
    \node[below left] at (0,0) {\tiny$O$};
    \draw[blue, thick] (-1.2,-0.8) -- (1.2,1.6);
    \draw[red, thick] (-1.2,-1.4) -- (1.2,1.0);
    \node[below, font=\scriptsize] at (0,-1.8) {无解（平行）};
    \node[right, font=\tiny, blue] at (0.8,1.4) {$k_1=k_2$};
    \node[right, font=\tiny, red] at (0.8,0.8) {$b_1{\neq}b_2$};
  \end{scope}

  % ====== 右：重合（无穷多解）======
  \begin{scope}[shift={(4.5, 0)}]
    \draw[->] (-1.5,0) -- (1.5,0) node[right] {\tiny$x$};
    \draw[->] (0,-1.5) -- (0,1.8) node[above] {\tiny$y$};
    \node[below left] at (0,0) {\tiny$O$};
    % 两线重合，画一条加粗的紫色线
    \draw[blue!60!red, very thick] (-1.2,-0.8) -- (1.2,1.6);
    \node[below, font=\scriptsize] at (0,-1.8) {无穷多解（重合）};
    \node[right, font=\tiny] at (0.6,1.2) {重合};
  \end{scope}

\end{tikzpicture}
\end{document}
